\section{AIME I 1990}
        \begin{enumerate}
        	\item (\textit{AIME I 1990}) The \href{https://artofproblemsolving.com/wiki/index.php?title=Increasing\_sequence}{increasing sequence} $2,3,5,6,7,10,11,\ldots$ consists of all \href{https://artofproblemsolving.com/wiki/index.php?title=Positive\_integer}{positive integers} that are neither the \href{https://artofproblemsolving.com/wiki/index.php?title=Perfect\_square}{ square} nor the \href{https://artofproblemsolving.com/wiki/index.php?title=Perfect\_cube}{ cube} of a positive integer. Find the 500th term of this sequence.


 
 \textbf{Answer: }528



\textbf{Solution 1}

Because there aren't that many perfect squares or cubes, let's look for the smallest perfect square greater than $500$. This happens to be $23^2=529$. Notice that there are $23$ squares and $8$ cubes less than or equal to $529$, but $1$ and $2^6$ are both squares and cubes. Thus, there are $529-23-8+2=500$ numbers in our sequence less than $529$.  Magically, we want the $500th$ term, so our answer is the biggest non-square and non-cube less than $529$, which is $\boxed{528}$.

 





\textbf{Solution 2}

This solution is similar as Solution 1, but to get the intuition why we chose to consider $23^2 = 529$, consider this:

 We need $n - T = 500$, where $n$ is an integer greater than 500 and $T$ is the set of numbers which contains all $k^2,k^3\le 500$.

 Firstly, we clearly need $n > 500$, so we substitute n for the smallest square or cube greater than $500$. However, if we use $n=8^3=512$, the number of terms in $T$ will exceed $n-500$. Therefore, $n=23^2=529$, and the number of terms in $T$ is $23+8-2=29$ by the \href{https://artofproblemsolving.com/wiki/index.php?title=Principle\_of\_Inclusion-Exclusion}{Principle of Inclusion-Exclusion}, fulfilling our original requirement of $n-T=500$.As a result, our answer is $529-1 = \boxed{528}$.

 




	\item (\textit{AIME I 1990}) Find the value of $(52+6\sqrt{43})^{3/2}-(52-6\sqrt{43})^{3/2}$.


 
 \textbf{Answer: }828



\textbf{Solution 1}

Suppose that $52+6\sqrt{43}$ is in the form of $(a + b\sqrt{43})^2$. \href{https://artofproblemsolving.com/wiki/index.php?title=FOIL}{FOILing} yields that $52 + 6\sqrt{43} = a^2 + 43b^2 + 2ab\sqrt{43}$. This implies that $a$ and $b$ equal one of $\pm3, \pm1$. The possible \href{https://artofproblemsolving.com/wiki/index.php?title=Set}{sets} are $(3,1)$ and $(-3,-1)$; the latter can be discarded since the \href{https://artofproblemsolving.com/wiki/index.php?title=Square\_root}{square root} must be positive. This means that $52 + 6\sqrt{43} = (\sqrt{43} + 3)^2$. Repeating this for $52-6\sqrt{43}$, the only feasible possibility is $(\sqrt{43} - 3)^2$.

 Rewriting, we get $(\sqrt{43} + 3)^3 - (\sqrt{43} - 3)^3$. Using the difference of \href{https://artofproblemsolving.com/wiki/index.php?title=Cube}{cubes}, we get that $[\sqrt{43} + 3\ - \sqrt{43} + 3]\ [(43 + 6\sqrt{43} + 9) + (43 - 9) + (43 - 6\sqrt{43} + 9)]$ $= (6)(3 \cdot 43 + 9) = \boxed{828}$.Note: You can also just use the formula $(a + b)^2 = a^2 + 2ab + b^2$ instead of \href{https://artofproblemsolving.com/wiki/index.php?title=FOIL}{ foiling}.

 





\textbf{Solution 2}

The $3/2$ power is quite irritating to work with so we look for a way to eliminate that. Notice that squaring the expression will accomplish that. Let $S$ be the sum of the given expression. 
 \[S^2= ((52+6\sqrt{43})^{3/2}-(52-6\sqrt{43})^{3/2})^2\]
 \[S^2 = (52+6\sqrt{43})^{3} + (52-6\sqrt{43})^{3} - 2((52+6\sqrt{43})(52-6\sqrt{43}))^{3/2}\]After doing the arithmetic (note that the first two terms will have some cancellation and that the last term will simplify quickly using difference of squares), we arrive at $S^2 = 685584$ which gives $S=\boxed{828}$.

 





\textbf{Solution 3}

Factor as a difference of cubes.
 \[\left[\left(52+6\sqrt{43}\right)^{\frac{1}{2}}-\left(52-6\sqrt{43}\right)^{\frac{1}{2}}\right]\left[\left(\left(\left(52+6\sqrt{43}\right)^{\frac{1}{2}}\right)^2+\left(52+6\sqrt{43}\right)^{\frac{1}{2}}\left(52-6\sqrt{43}\right)^{\frac{1}{2}}+\left(\left(52-6\sqrt{43}\right)^{\frac{1}{2}}\right)^2\right)\right] =\]
 \[\left[\left(52+6\sqrt{43}\right)^{\frac{1}{2}}-\left(52-6\sqrt{43}\right)^{\frac{1}{2}}\right]\left[104+\left(52^2-\left(36\right)\left(43\right)\right)^{\frac{1}{2}}\right] =\]
 \[\left[\left(52+6\sqrt{43}\right)^{\frac{1}{2}}-\left(52-6\sqrt{43}\right)^{\frac{1}{2}}\right]\left[104+34\right].\]We can simplify the left factor as follows.
 \[\left(52+6\sqrt{43}\right)^{\frac{1}{2}}-\left(52-6\sqrt{43}\right)^{\frac{1}{2}} = x\]
 \[104-2\left(52+6\sqrt{43}\right)^{\frac{1}{2}}\left(52-6\sqrt{43}\right)^{\frac{1}{2}} = x^2\]
 \[104-68 = x^2\]
 \[36 = x^2.\]Since $\left(52+6\sqrt{43}\right)^{\frac{1}{2}} > \left(52-6\sqrt{43}\right)^{\frac{1}{2}}$, we know that $x=6$, so our final answer is $(6)(138) = \boxed{828}$.

 


 





\textbf{Solution 4}

Let $x=52+6\sqrt{43}$, $y=52-6\sqrt{43}$. Similarly to solution 2, we let 
 \[S=x^{\frac{3}{2}}+y^{\frac{3}{2}}\]
 \begin{align*} S^2&=(x^{\frac{3}{2}}+y^{\frac{3}{2}})^2\\ &=x^3+y^3+2x^{\frac{3}{2}}y^{\frac{3}{2}} \end{align*}
The expression can be simplified as follow
 \begin{align*} S^2&=x^3+y^3+2x^{\frac{3}{2}}y^{\frac{3}{2}}\\ &=(x+y)(x^2-xy+y^2)+2(xy)^{\frac{3}{2}}\\ &=(x+y)((x+y)^2-xy)+2\sqrt{xy}^3\\ &=(x+y)((x+y)^2-\sqrt{xy}^2)+\sqrt{xy}^3\\ &=(x+y)(x+y+\sqrt{xy})(x+y-\sqrt{xy})+2\sqrt{xy}^3\\ &=104(104+34)(104-34)+2\cdot34^3\\ &=685584 \end{align*}
Thus $S=\sqrt{685584}=\boxed{828}$.

 I don't believe we get 685584 from 104(138)(70) + 2*34\textasciicircum{}3 though?

 





\textbf{Solution 5}

(Similar to Solution 3, but with substitution)

 Let $a=\sqrt{52+6\sqrt{43}}$ and $b=\sqrt{52-6\sqrt{43}}.$ We want to find $a^3-b^3=(a-b)(a^2+ab+b^2).$

 We have
 \[a^2+b^2=104,\text{ and}\]
 \[ab=\sqrt{(52+6\sqrt{43})(52-6\sqrt{43})}=\sqrt{1156}=34.\]Then, $(a-b)^2=a^2+b^2-2ab=104-2\cdot 34= 36\implies a-b=6.$

 Our answer is
 \[a^3-b^3=(a-b)(a^2+b^2+ab)=6\cdot 138=\boxed{828.}\]

 





\textbf{Solution 6}

(Similar to Solution 1, but expanding the cubes instead)

 Like in Solution 1, we have $\sqrt{52 + 6\sqrt{43}} = \sqrt{43} + 3$ and $\sqrt{52 - 6\sqrt{43}} = \sqrt{43} - 3.$ 

 Therefore we have that $(52 + 6\sqrt{43})^{3/2} - (52 + 6\sqrt{43})^{3/2}$ $= \sqrt{52 + 6\sqrt{43}}^3 - \sqrt{52 - 6\sqrt{43}}^3$ $= (\sqrt{43} + 3)^3 - (\sqrt{43} - 3)^3.$

 From here, we use the formula $(a + b)^3 = a^3 + 3a^2b + 3ab^2 + b^3$ and $(a - b)^3 = a^3 - 3a^2b + 3ab^2 - b^3$. Applying them to our problem we get that $(\sqrt{43} + 3)^3 - (\sqrt{43} - 3)^3 = (27 + 27\sqrt{43} + 9 \cdot 43 + 43\sqrt{43}) - (-27 + 27\sqrt{43} - 9*43 + 43\sqrt{43}).$ We see that all the terms with square roots cancel, leaving us with $2 (27 + 9 \cdot 43) = 2 \cdot 414 = \boxed{828}.$

 Note: We have that $\sqrt{52 - 6\sqrt{43}} = \sqrt{43} - 3$ because we need the square root to be positive and $\sqrt{43} > 3$ since $43$ is obviously greater than $9.$ So we have $\sqrt{52 - 6\sqrt{43}} = \sqrt{43} - 3.$

 




	\item (\textit{AIME I 1990}) Let $P_1^{}$ be a regular $r~\mbox{gon}$ and $P_2^{}$ be a regular $s~\mbox{gon}$ $(r\geq s\geq 3)$ such that each interior angle of $P_1^{}$ is $\frac{59}{58}$ as large as each interior angle of $P_2^{}$. What's the largest possible value of $s_{}^{}$?


 
 \textbf{Answer: }117



\textbf{Solution 1}

The formula for the interior angle of a regular sided \href{https://artofproblemsolving.com/wiki/index.php?title=Polygon}{polygon} is $\frac{(n-2)180}{n}$. 

 Thus, $\frac{\frac{(r-2)180}{r}}{\frac{(s-2)180}{s}} = \frac{59}{58}$. Cross multiplying and simplifying, we get $\frac{58(r-2)}{r} = \frac{59(s-2)}{s}$. Cross multiply and combine like terms again to yield $58rs - 58 \cdot 2s = 59rs - 59 \cdot 2r \Longrightarrow 118r - 116s = rs$. Solving for $r$, we get $r = \frac{116s}{118 - s}$.

 $r \ge 0$ and $s \ge 0$, making the \href{https://artofproblemsolving.com/wiki/index.php?title=Numerator}{numerator} of the \href{https://artofproblemsolving.com/wiki/index.php?title=Fraction}{fraction} positive. To make the \href{https://artofproblemsolving.com/wiki/index.php?title=Denominator}{denominator} \href{https://artofproblemsolving.com/wiki/index.php?title=Positive}{positive}, $s < 118$; the largest possible value of $s$ is $117$.

 This is achievable because the denominator is $1$, making $r$ a positive number $116 \cdot 117$ and $s = \boxed{117}$.

 





\textbf{Solution 2}

Like above, use the formula for the interior angles of a regular sided \href{https://artofproblemsolving.com/wiki/index.php?title=Polygon}{polygon}.

 
$\frac{(r-2)180}{r} = \frac{59}{58} * \frac{(s-2)180}{s}$

 
$59 * 180 * (s-2) * r = 58 * 180 * (r-2) * s$

 
$59 * (rs - 2r) = 58 * (rs - 2s)$

 
$rs - 118r = -116s$

 
$rs = 118r-116s$

 
This equation tells us $s$ divides $118r$. If $s$ specifically divides 118 then the highest it can be is 118. However, this gives an equation with no solution. The second largest possibility in this case is $s=59$, which does give a solution: $s=59, r=116$. Although, the problem asks for $s$, not $r$. The only conceivable reasoning behind this is that $r$ is greater than 1000. This prompts us to look into the second case, where $s$ divides $r$. Make $r = s * k$. Rewrite the equation using this new information.

 
$s * s * k = 118 * s * k - 116 * s$

 
$s * k = 118 * k - 116$

 
Now we now k divides 116. The larger k is, the larger s will be, so we set k to be the maximum: 116.

 
$s * 116 = 118 * 116 - 116$

 
$s = 118 - 1$

 
$s = \boxed{117}$

 
-jackshi2006

 


 





\textbf{Solution 3}

As in above, we have $rs = 118r - 116s.$ This means that $rs + 116s - 118r = 0.$ Using SFFT we obtain $s(r+116) - 118(r+116) = -118 \cdot 116 \implies (s-118)(r+116) = -118 \cdot 116.$ Since $r+116$ is always positive, we know that $s-118$ must be negative. Therefore the maximum value of $s$ must be $\boxed{117}$ which indeed yields an integral value of $r.$

 




	\item (\textit{AIME I 1990}) Find the positive solution to


 $\frac 1{x^2-10x-29}+\frac1{x^2-10x-45}-\frac 2{x^2-10x-69}=0$
 \textbf{Answer: }013



\textbf{Solution}

We could clear out the denominators by multiplying, though that would be unnecessarily tedious.

 To simplify the equation, substitute $a = x^2 - 10x - 29$ (the denominator of the first fraction). We can rewrite the equation as $\frac{1}{a} + \frac{1}{a - 16} - \frac{2}{a - 40} = 0$. Multiplying out the denominators now, we get:

 
 \[(a - 16)(a - 40) + a(a - 40) - 2(a)(a - 16) = 0\]

 Simplifying, $-64a + 40 \times 16 = 0$, so $a = 10$. Re-substituting, $10 = x^2 - 10x - 29 \Longleftrightarrow 0 = (x - 13)(x + 3)$. The positive \href{https://artofproblemsolving.com/wiki/index.php?title=Root}{root} is $\boxed{013}$.

 




	\item (\textit{AIME I 1990}) Let $n^{}_{}$ be the smallest positive integer that is a multiple of $75_{}^{}$ and has exactly $75_{}^{}$ positive integral divisors, including $1_{}^{}$ and itself. Find $\frac{n}{75}$.


 
 \textbf{Answer: }432



\textbf{Solution}

The \href{https://artofproblemsolving.com/wiki/index.php?title=Prime\_factorization}{prime factorization} of $75 = 3^15^2 = (2+1)(4+1)(4+1)$. For $n$ to have exactly $75$ integral divisors, we need to have $n = p_1^{e_1-1}p_2^{e_2-1}\cdots$ such that $e_1e_2 \cdots = 75$. Since $75|n$, two of the \href{https://artofproblemsolving.com/wiki/index.php?title=Prime}{prime} \href{https://artofproblemsolving.com/wiki/index.php?title=Factor}{factors} must be $3$ and $5$. To minimize $n$, we can introduce a third prime factor, $2$. Also to minimize $n$, we want $5$, the greatest of all the factors, to be raised to the least power. Therefore, $n = 2^43^45^2$ and $\frac{n}{75} = \frac{2^43^45^2}{3 \cdot 5^2} = 16 \cdot 27 = \boxed{432}$.

 $\textbf{ Note: the smallest integer n which has 75 factors, happens to be a multiple of 75 anyway  }$

 Note: the problem where this is linked is covered 0:40-1:40

 




	\item (\textit{AIME I 1990}) A biologist wants to calculate the number of fish in a lake. On May 1 she catches a random sample of 60 fish, tags them, and releases them. On September 1 she catches a random sample of 70 fish and finds that 3 of them are tagged. To calculate the number of fish in the lake on May 1, she assumes that 25\% of these fish are no longer in the lake on September 1 (because of death and emigrations), that 40\% of the fish were not in the lake May 1 (because of births and immigrations), and that the number of untagged fish and tagged fish in the September 1 sample are representative of the total population. What does the biologist calculate for the number of fish in the lake on May 1? 


 
 \textbf{Answer: }840



\textbf{Solution 1}

Of the $70$ fish caught in September, $40\%$ were not there in May, so $42$ fish were there in May. Since the percentage of tagged fish in September is proportional to the percentage of tagged fish in May, $\frac{3}{42} = \frac{60}{x} \Longrightarrow \boxed{x = 840}$.

 (Note the 25\% death rate does not affect the answer because both tagged and nontagged fish die.)

 





\textbf{Solution 2}

First, we notice that there are 45 tags left, after 25\% of the original fish have went away/died. Then, some $x$ percent of fish have been added such that $\frac{x}{x+75} = 40 \%$, or $\frac{2}{5}$. Solving for $x$, we get that $x = 50$, so the total number of fish in September is $125 \%$, or $\frac{5}{4}$ times the total number of fish in May.

 Since $\frac{3}{70}$ of the fish in September were tagged, $\frac{45}{5n/4} = \frac{3}{70}$, where $n$ is the number of fish in May. Solving for $n$, we see that $n = \boxed{840}$

 





\textbf{Solution 3(easy to comprehend)}

Let $k$ be the number of fish in the lake on May 1, and $n$ be the number of fish in the lake on September 1. Because the biologist believes that 40\% of $n$ were not there on May 1, we have the equality 
 \[\frac{6}{10}\cdot n = \frac{3}{4} \cdot k\] which reduces to 
 \[n=\frac{5}{4} \cdot k\] Then, since we have 3 out of 70 fishies tagged on September 1st, we can write the equality 
 \[\frac{3}{70}=\frac{60\cdot\frac{3}{4}}{\frac{5\cdot k}{4}}\] because 25\% of the 60 fishes tagged on May 1st are now gone, for the numerator, and for the deonominator we have $n$, which we found is $\frac{5}{4}\cdot k$. Solving, we get $k= \boxed{840}$

 




	\item (\textit{AIME I 1990}) A triangle has vertices $P_{}^{}=(-8,5)$, $Q_{}^{}=(-15,-19)$, and $R_{}^{}=(1,-7)$. The equation of the bisector of $\angle P$ can be written in the form $ax+2y+c=0_{}^{}$. Find $a+c_{}^{}$.


 
 \textbf{Answer: }089



\textbf{Solution}

Use the \href{https://artofproblemsolving.com/wiki/index.php?title=Distance\_formula}{distance formula} to determine the lengths of each of the sides of the triangle. We find that it has lengths of side $15,\ 20,\ 25$, indicating that it is a $3-4-5$ \href{https://artofproblemsolving.com/wiki/index.php?title=Right\_triangle}{right triangle}. At this point, we just need to find another \href{https://artofproblemsolving.com/wiki/index.php?title=Point}{point} that lies on the bisector of $\angle P$.

 





\textbf{Solution 1}


                    \begin{figure}[H]
                        \centering
                        \includegraphics[width=0.5\textwidth]{Solution/AIME_1990_I_7_Solution_Figure_1.png}
                    \end{figure}
                    Use the \href{https://artofproblemsolving.com/wiki/index.php?title=Angle\_bisector\_theorem}{angle bisector theorem} to find that the angle bisector of $\angle P$ divides $QR$ into segments of length $\frac{25}{x} = \frac{15}{20 -x} \Longrightarrow x = \frac{25}{2},\ \frac{15}{2}$. It follows that $\frac{QP'}{RP'} = \frac{5}{3}$, and so $P' = \left(\frac{5x_R + 3x_Q}{8},\frac{5y_R + 3y_Q}{8}\right) = (-5,-23/2)$.

 The desired answer is the equation of the line $PP'$. $PP'$ has slope $\frac{-11}{2}$, from which we find the equation to be $11x + 2y + 78 = 0$. Therefore, $a+c = \boxed{089}$. 

 





\textbf{Solution 2}


                    \begin{figure}[H]
                        \centering
                        \includegraphics[width=0.5\textwidth]{Solution/AIME_1990_I_7_Solution_Figure_1.png}
                    \end{figure}
                    Extend $PR$ to a point $S$ such that $PS = 25$. This forms an \href{https://artofproblemsolving.com/wiki/index.php?title=Isosceles\_triangle}{isosceles triangle} $PQS$. The \href{https://artofproblemsolving.com/wiki/index.php?title=Coordinate}{coordinates} of $S$, using the slope of $PR$ (which is $-4/3$), can be determined to be $(7,-15)$. Since the \href{https://artofproblemsolving.com/wiki/index.php?title=Angle\_bisector}{angle bisector} of $\angle P$ must touch the midpoint of $QS \Rightarrow (-4,-17)$, we have found our two points. We reach the same answer of $11x + 2y + 78 = 0$.

 





\textbf{Solution 3}


                    \begin{figure}[H]
                        \centering
                        \includegraphics[width=0.5\textwidth]{Solution/AIME_1990_I_7_Solution_Figure_1.png}
                    \end{figure}
                    By the angle bisector theorem as in solution 1, we find that $QP' = 25/2$. If we draw the right triangle formed by $Q, P',$ and the point directly to the right of $Q$ and below $P'$, we get another $3-4-5 \triangle$ (since the slope of $QR$ is $3/4$). Using this, we find that the horizontal projection of $QP'$ is $10$ and the vertical projection of $QP'$ is $15/2$. 

 Thus, the angle bisector touches $QR$ at the point $\left(-15 + 10, -19 + \frac{15}{2}\right) = \left(-5,-\frac{23}{2}\right)$, from where we continue with the first solution.

 





\textbf{Solution 4}

This solution uses terminology from the other solutions. The incenter is a much easier point to find on the line $PP'$. Note that the inradius of $\triangle PQR$ is $5$. If you do not understand this, substitute values into the $[\triangle ABC] = rs$ equation. If lines are drawn from the incenter perpendicular to $PR$ and $QR$, then a square with side length $5$ will be created. Call the point opposite $R$ in this square $R'$. Since $R$ has coordinates $(1, -7)$, and the sides of the squares are on a $3-4-5$ ratio, the coordinates of $R'$ are $(-6, -6)$. This is because the x-coordinate is moving to the left $4+3=7$ units and the y-coordinate is moving up $-3+4=1$ units. The line through $(-8,5)$ and $(-6,-6)$ is $11x+2y+78=0$.

 





\textbf{Solution 5 (Trigonometry)}


                    \begin{figure}[H]
                        \centering
                        \includegraphics[width=0.5\textwidth]{Solution/AIME_1990_I_7_Solution_Figure_1.png}
                    \end{figure}
                    Transform triangle $PQR$ so that $P$ is at the origin. Note that the slopes do not change when we transform the triangle.

 
We know that the slope of $PQ$ is $\frac{24}{7}$ and the slope of $PR$ is $-\frac{4}{3}$. Thus, in the complex plane, they are equivalent to $\tan(\alpha)=\frac{24}{7}$ and $\tan(\beta)=-\frac{4}{3}$, respectively. Here $\alpha$ is the angle formed by the $x$-axis and $PQ$, and $\beta$ is the angle formed by the $x$-axis and $PR$. The equation of the angle bisector is $\tan\left(\frac{\alpha+\beta}{2}\right)$.

 
As the tangents are in very neat \href{https://artofproblemsolving.com/wiki/index.php?title=Pythagorean\_triple}{Pythagorean triples}, we can easily calculate $\cos(\alpha)$ and $\cos(\beta)$.

 Angle $\alpha$ is in the third quadrant, so $\cos(\alpha)$ is negative. Thus $\cos(\alpha)=-\frac{7}{25}$.

 Angle $\beta$ is in the fourth quadrant, so $\cos(\beta)$ is positive. Thus $\cos(\beta)=\frac{3}{5}$.

 
By the \href{https://artofproblemsolving.com/wiki/index.php?title=Trigonometric\_identities\#Half-angle\_identities}{Half-Angle Identities}, $\tan\left(\frac{\alpha}{2}\right)=\pm\sqrt{\frac{1-\cos(\alpha)}{1+\cos(\alpha)}}=\pm\sqrt{\frac{\frac{32}{25}}{\frac{18}{25}}}=\pm\sqrt{\frac{16}{9}}=\pm\frac{4}{3}$ and $\tan\left(\frac{\beta}{2}\right)=\pm\sqrt{\frac{1-\cos(\beta)}{1+\cos(\beta)}}=\pm\sqrt{\frac{\frac{2}{5}}{\frac{8}{5}}}=\pm\sqrt{\frac{1}{4}}=\pm\frac{1}{2}$.

 Since $\frac{\alpha}{2}$ and $\frac{\beta}{2}$ must be in the second quadrant, their tangent values are both negative. Thus $\tan\left(\frac{\alpha}{2}\right)=-\frac{4}{3}$ and $\tan\left(\frac{\beta}{2}\right)=-\frac{1}{2}$.

 
By the \href{https://artofproblemsolving.com/wiki/index.php?title=Trigonometric\_identities\#Angle\_addition\_identities}{sum of tangents formula}, $\tan\left(\frac{\alpha}{2}+\frac{\beta}{2}\right)=\frac{\tan\left(\frac{\alpha}{2}\right)+\tan\left(\frac{\beta}{2}\right)}{1-\tan\left(\frac{\alpha}{2}\right)\tan\left(\frac{\beta}{2}\right)}=\frac{-\frac{11}{6}}{\frac{1}{3}}=-\frac{11}{2}$, which is the slope of the angle bisector.

 
Finally, the equation of the angle bisector is $y-5=-\frac{11}{2}(x+8)$ or $y=-\frac{11}{2}x-39$. Rearranging, we get $11x+2y+78=0$, so our sum is $a+c=11+78=\boxed{089}$. ~eevee9406

 




	\item (\textit{AIME I 1990}) In a shooting match, eight clay targets are arranged in two hanging columns of three targets each and one column of two targets. A marksman is to break all the targets according to the following rules: 


 1) The marksman first chooses a column from which a target is to be broken. 


 2) The marksman must then break the lowest remaining target in the chosen column. 


 If the rules are followed, in how many different orders can the eight targets be broken?


 
                    \begin{figure}[H]
                        \centering
                        \includegraphics[width=0.5\textwidth]{Problem/AIME_1990_I_8_Problem_Figure_1.png}
                    \end{figure}
                    
 \textbf{Answer: }560



\textbf{Solution}

Clearly, the marksman must shoot the left column three times, the middle column two times, and the right column three times.

 From left to right, suppose that the columns are labeled $L,M,$ and $R,$ respectively. We consider the string $LLLMMRRR:$

 Since the letter arrangements of $LLLMMRRR$ and the shooting orders have one-to-one correspondence, we count the letter arrangements: 
 \[\binom{8}{3,2,3} = \frac{8!}{3!\cdot2!\cdot3!} = \boxed{560}.\]

 \textbf{Remark} We can count the letter arrangements of $LLLMMRRR$ in two ways:

 <ol style="margin-left: 1.5em;"> <li>There are $8!$ ways to arrange $8$ distinguishable letters. However: 

 <ul style="list-style-type:square;">\item Since there are $3$ indistinguishable $L$'s, we have counted each distinct arrangement of the $L$'s for $3!$ times. 

 \item Since there are $2$ indistinguishable $M$'s, we have counted each distinct arrangement of the $M$'s for $2!$ times. 

 \item Since there are $3$ indistinguishable $R$'s, we have counted each distinct arrangement of the $R$'s for $3!$ times. 

 </ul>By the Multiplication Principle, we have counted each distinct arrangement of $LLLMMRRR$ for $3!\cdot2!\cdot3!$ times. As shown in the \textbf{Solution} section, we use division to fix the overcount. The answer is $\frac{8!}{3!\cdot2!\cdot3!} = 560.$ 

 Alternatively, we can use a multinomial coefficient to obtain the answer directly: $\binom{8}{3,2,3}=\frac{8!}{3!\cdot2!\cdot3!} = 560.$ 

 </li> 

 <li>First, we have $\binom83$ ways to choose any $3$ from the $8$ positions for the $L$'s.Next, we have $\binom52$ ways to choose any $2$ from the $5$ remaining positions for the $M$'s.

 Finally, we have $\binom33$ way to choose $3$ from the $3$ remaining positions for the $R$'s.

 By the Multiplication Principle, the answer is $\binom83\binom52\binom33=560.$

 </li>

 </ol>




	\item (\textit{AIME I 1990}) A fair coin is to be tossed $10_{}^{}$ times. Let $i/j^{}_{}$, in lowest terms, be the probability that heads never occur on consecutive tosses. Find $i+j_{}^{}$. 


 
 \textbf{Answer: }073



\textbf{Solution 1}

Clearly, at least $5$ tails must be flipped; any less, then by the \href{https://artofproblemsolving.com/wiki/index.php?title=Pigeonhole\_Principle}{Pigeonhole Principle} there will be heads that appear on consecutive tosses. 

 Consider the case when $5$ tails occur. The heads must fall between the tails such that no two heads fall between the same tails, and must fall in the positions labeled $(H)$:

 
            \begin{enumerate}[label=(\roman*)]
                \item $(H)\ T\ (H)\ T\ (H)\ T\ (H)\ T\ (H)\ T\ (H)$
            \end{enumerate}
            There are six slots for the heads to be placed, but only $5$ heads remaining. Thus, using \href{https://artofproblemsolving.com/wiki/index.php?title=Stars-and-bars}{stars-and-bars} there are ${6\choose5}$ possible combinations of $5$ heads. Continuing this pattern, we find that there are $\sum_{i=6}^{11} {i\choose{11-i}} = {6\choose5} + {7\choose4} + {8\choose3} + {9\choose2} + {{10}\choose1} + {{11}\choose0} = 144$. There are a total of $2^{10}$ possible flips of $10$ coins, making the probability $\frac{144}{1024} = \frac{9}{64}$. Thus, our solution is $9 + 64 = \boxed{073}$.

 





\textbf{Solution 2 (Recursion)}

Call the number of ways of flipping $n$ coins and not receiving any consecutive heads $S_n$. Notice that tails must be received in at least one of the first two flips.

 If the first coin flipped is a T, then the remaining $n-1$ flips must fall under one of the configurations of $S_{n-1}$.

 If the first coin flipped is a H, then the second coin must be a T. There are then $S_{n-2}$ configurations.

 Thus, $S_n = S_{n-1} + S_{n-2}$. By counting, we can establish that $S_1 = 2$ and $S_2 = 3$. Therefore, $S_3 = 5,\ S_4 = 8$, forming the \href{https://artofproblemsolving.com/wiki/index.php?title=Fibonacci\_sequence}{Fibonacci sequence}. Listing them out, we get $2,3,5,8,13,21,34,55,89,144$, and the 10th number is $144$. Putting this over $2^{10}$ to find the probability, we get $\frac{9}{64}$. Our solution is $9+64=\boxed{073}$.

 





\textbf{Solution 3}

We can also split the problem into casework.

 Case 1: 0 Heads

 There is only one possibility.

 Case 2: 1 Head

 There are 10 possibilities.

 Case 3: 2 Heads

 There are 36 possibilities.

 Case 4: 3 Heads

 There are 56 possibilities.

 Case 5: 4 Heads

 There are 35 possibilities.

 Case 6: 5 Heads

 There are 6 possibilities.

 We have $1+10+36+56+35+6=144$, and there are $1024$ possible outcomes, so the probability is $\frac{144}{1024}=\frac{9}{64}$, and the answer is $\boxed{073}$.

 $\textbf{-RootThreeOverTwo}$

 




	\item (\textit{AIME I 1990}) The sets $A = \{z : z^{18} = 1\}$ and $B = \{w : w^{48} = 1\}$ are both sets of complex roots of unity.  The set $C = \{zw : z \in A ~ \mbox{and} ~ w \in B\}$ is also a set of complex roots of unity.  How many distinct elements are in $C^{}_{}$?


 
 \textbf{Answer: }144



\textbf{Solution 1}

The \href{https://artofproblemsolving.com/wiki/index.php?title=Least\_common\_multiple}{least common multiple} of $18$ and $48$ is $144$, so define $n = e^{2\pi i/144}$. We can write the numbers of set $A$ as $\{n^8, n^{16}, \ldots n^{144}\}$ and of set $B$ as $\{n^3, n^6, \ldots n^{144}\}$. $n^x$ can yield at most $144$ different values. All solutions for $zw$ will be in the form of $n^{8k_1 + 3k_2}$. 

 $8$ and $3$ are relatively prime, and by the \href{https://artofproblemsolving.com/wiki/index.php?title=Chicken\_McNugget\_Theorem}{Chicken McNugget Theorem}, for two relatively prime integers $a,b$, the largest number that cannot be expressed as the sum of multiples of $a,b$ is $a \cdot b - a - b$. For $3,8$, this is $13$; however, we can easily see that the numbers $145$ to $157$ can be written in terms of $3,8$. Since the exponents are of roots of unities, they reduce $\mod{144}$, so all numbers in the range are covered. Thus the answer is $\boxed{144}$.

 \textbf{Comment of S1} Using \href{https://artofproblemsolving.com/wiki/index.php?title=Chicken\_McNugget\_Theorem}{Chicken McNugget Theorem} and the property of complex numbers for reasoning is sharp, but note that the Frobenius Number is only constrained to positive integer pairs. Please check on "Comment of S2" below to see how to use \href{https://artofproblemsolving.com/wiki/index.php?title=Diophantine\_equation}{Diophantine equation} to make a simple deduction. ~ Will\_Dai

 





\textbf{Solution 2}

The 18 and 48th roots of $1$ can be found by \href{https://artofproblemsolving.com/wiki/index.php?title=De\_Moivre\%27s\_Theorem}{De Moivre's Theorem}. They are $\text{cis}\,\left(\frac{2\pi k_1}{18}\right)$ and $\text{cis}\,\left(\frac{2\pi k_2}{48}\right)$ respectively, where $\text{cis}\,\theta = \cos \theta + i \sin \theta$ and $k_1$ and $k_2$ are integers from $0$ to $17$ and $0$ to $47$, respectively. 

 $zw = \text{cis}\,\left(\frac{\pi k_1}{9} + \frac{\pi k_2}{24}\right) = \text{cis}\,\left(\frac{8\pi k_1 + 3 \pi k_2}{72}\right)$. Since the \href{https://artofproblemsolving.com/wiki/index.php?title=Trigonometry}{trigonometric} functions are \href{https://artofproblemsolving.com/wiki/index.php?title=Periodic\_function}{periodic} every $2\pi$, there are at most $72 \cdot 2 = \boxed{144}$ distinct elements in $C$. As above, all of these will work.

 \textbf{Comment on S2} By the property of \href{https://artofproblemsolving.com/wiki/index.php?title=Diophantine\_equation}{Diophantine equation}, given a problem to find integers $x$ and $y$ so that $ax + by = c$ for some integer constants $a$, $b$, $c$: if $\gcd(a, b) = 1$, then any arbitrary integer $c$ could by formed with some combination of integers $(x, y)$. Double checking the constraints of $k_1$ and $k_2$, we realize that all integers of $[0, 143]$ can be formed by $8k_1 + 3k_2$, yielding the answer of $144$.~Will\_Dai

 





\textbf{Solution 3}

The values in polar form will be $(1, 20x)$ and $(1, 7.5x)$. Multiplying these gives $(1, 27.5x)$. Then, we get $27.5$, $55$, $82.5$, $110$, $...$ up to $3960$ $(\text{lcm}(55,360)) \implies \frac{3960 \cdot 2}{55}=\boxed{144}$.

 




	\item (\textit{AIME I 1990}) Someone observed that $6! = 8 \cdot 9 \cdot 10$.  Find the largest positive integer $n^{}_{}$ for which $n^{}_{}!$ can be expressed as the product of $n - 3_{}^{}$ consecutive positive integers.


 
 \textbf{Answer: }023



\textbf{Solution 1}

The product of $n - 3$ consecutive integers can be written as $\frac{(n - 3 + a)!}{a!}$ for some integer $a$. Thus, $n! = \frac{(n - 3 + a)!}{a!}$, from which it  becomes evident that $a \ge 3$. Since $(n - 3 + a)! > n!$, we can rewrite this as $\frac{n!(n+1)(n+2) \ldots (n-3+a)}{a!} = n! \Longrightarrow (n+1)(n+2) \ldots (n-3+a) = a!$. For $a = 4$, we get $n + 1 = 4!$ so $n = 23$. For greater values of $a$, we need to find the product of $a-3$ consecutive integers that equals $a!$. $n$ can be approximated as $^{a-3}\sqrt{a!}$, which decreases as $a$ increases. Thus, $n = 23$ is the greatest possible value to satisfy the given conditions.

 





\textbf{Solution 2}

Let the largest of the $n-3$ consecutive positive integers be $k$. Clearly $k$ cannot be less than or equal to $n$, else the product of $n-3$ consecutive positive integers will be less than $n!$.

 Key observation:Now for $n$ to be maximum the smallest number (or starting number) of the $n-3$ consecutive positive integers must be minimum, implying that $k$ needs to be minimum. But the least $k > n$ is $n+1$.

 So the $n-3$ consecutive positive integers are $5, 6, 7\ldots, n+1$

 So we have $\frac{(n+1)!}{4!} = n!$$\Longrightarrow  n+1 = 24$$\Longrightarrow  n = 23$

 \textbf{Generalization} Largest positive integer $n$ for which $n!$ can be expressed as the product of $n-a$ consecutive positive integers is $(a+1)! - 1$

 For ex. largest $n$ such that product of $n-6$ consecutive positive integers is equal to $n!$ is $7!-1 = 5039$

 Proof:Reasoning the same way as above, let the largest of the $n-a$ consecutive positive integers be $k$. Clearly $k$ cannot be less than or equal to $n$, else the product of $n-a$ consecutive positive integers will be less than $n!$.

 Now, observe that for $n$ to be maximum the smallest number (or starting number) of the $n-a$ consecutive positive integers must be minimum, implying that $k$ needs to be minimum. But the least $k > n$ is $n+1$.

 So the $n-a$ consecutive positive integers are $a+2, a+3, \ldots n+1$

 So we have $\frac{(n+1)!}{(a+1)!} = n!$$\Longrightarrow  n+1 = (a+1)!$$\Longrightarrow  n = (a+1)! -1$

 $Kris17$

 




	\item (\textit{AIME I 1990}) A regular 12-gon is inscribed in a circle of radius 12.  The sum of the lengths of all sides and diagonals of the 12-gon can be written in the form


 $a + b \sqrt{2} + c \sqrt{3} + d \sqrt{6},$where $a^{}_{}$, $b^{}_{}$, $c^{}_{}$, and $d^{}_{}$ are positive integers.  Find $a + b + c + d^{}_{}$.


 
 \textbf{Answer: }720



\textbf{Solution 1}


                    \begin{figure}[H]
                        \centering
                        \includegraphics[width=0.5\textwidth]{Solution/AIME_1990_I_12_Solution_Figure_1.png}
                    \end{figure}
                    

 The easiest way to do this seems to be to find the length of each of the sides and diagonals. To do such, draw the \href{https://artofproblemsolving.com/wiki/index.php?title=Radius}{radii} that meet the endpoints of the sides/diagonals; this will form \href{https://artofproblemsolving.com/wiki/index.php?title=Isosceles\_triangle}{isosceles triangles}. Drawing the \href{https://artofproblemsolving.com/wiki/index.php?title=Altitude}{altitude} of those \href{https://artofproblemsolving.com/wiki/index.php?title=Triangle}{triangles} and then solving will yield the respective lengths.

 
                \begin{itemize}
                    \item The length of each of the 12 sides is $2 \cdot 12\sin 15$. $24\sin 15 = 24\sin (45 - 30) = 24\frac{\sqrt{6} - \sqrt{2}}{4} = 6(\sqrt{6} - \sqrt{2})$.\item The length of each of the 12 diagonals that span across 2 edges is $2 \cdot 12\sin 30 = 12$ (or notice that the triangle formed is \href{https://artofproblemsolving.com/wiki/index.php?title=Equilateral\_triangle}{equilateral}).\item The length of each of the 12 diagonals that span across 3 edges is $2 \cdot 12\sin 45 = 12\sqrt{2}$ (or notice that the triangle formed is a $45 - 45 - 90$ \href{https://artofproblemsolving.com/wiki/index.php?title=Right\_triangle}{right triangle}).\item The length of each of the 12 diagonals that span across 4 edges is $2 \cdot 12\sin 60 = 12\sqrt{3}$.\item The length of each of the 12 diagonals that span across 5 edges is $2 \cdot 12\sin 75 = 24\sin (45 + 30) = 24\frac{\sqrt{6}+\sqrt{2}}{4} = 6(\sqrt{6}+\sqrt{2})$.\item The length of each of the 6 \href{https://artofproblemsolving.com/wiki/index.php?title=Diameter}{diameters} is $2 \cdot 12 = 24$.
                \end{itemize}


                Adding all of these up, we get $12[6(\sqrt{6} - \sqrt{2}) + 12 + 12\sqrt{2} + 12\sqrt{3} + 6(\sqrt{6}+\sqrt{2})] + 6 \cdot 24$ 

 $= 12(12 + 12\sqrt{2} + 12\sqrt{3} + 12\sqrt{6}) + 144 = 288 + 144\sqrt{2} + 144\sqrt{3} + 144\sqrt{6}$. Thus, the answer is $144 \cdot 5 = \boxed{720}$.

 





\textbf{Solution 2}

A second method involves drawing a triangle connecting the center of the 12-gon to two vertices of the 12-gon.  Since the distance from the center to a vertex of the 12-gon is $12$, the Law of Cosines can be applied to this isosceles triangle, to give:

 $a^2 = 12^2 + 12^2 - 2\cdot 12\cdot 12\cdot \cos \theta$ 

 $a^2 = 2\cdot 12^2 - 2\cdot 12^2 \cos \theta$

 $a^2 = 2\cdot 12^2 (1 -  \cos \theta)$

 $a = 12\sqrt{2} \cdot \sqrt{1 - \cos \theta}$

 There are six lengths of sides/diagonals, corresponding to $\theta = {30^{\circ}, 60^{\circ}, 90^{\circ}, 120^{\circ}, 150^{\circ}, 180^{\circ}}$

 Call these lengths $a_1, a_2, a_3, a_4, a_5, a_6$ from shortest to longest.  The total length $l$ that is asked for is

 $l = 12(a_1 + a_2 + a_3 + a_4 + a_5) + 6a_6$, noting that $a_6$ as written gives the diameter of the circle, which is the longest diagonal.

 


 $l = 12[12\sqrt{2} (\sqrt{1 - \cos 30^{\circ}}+\sqrt{1-\cos 60^{\circ}}+\sqrt{1-\cos 90^{\circ}}+\sqrt{1-\cos 120^{\circ}}+\sqrt{1 - \cos 150^{\circ}})] + 6 \cdot 24$ 

 
$l = 144\sqrt{2} \left(\sqrt{1 - \frac{\sqrt{3}}{2}} + \sqrt{1 - \frac{1}{2}} + \sqrt{1-0} + \sqrt{1 + \frac{1}{2}} + \sqrt{1 + \frac{\sqrt{3}}{2}}\right) + 144$

 
$l = 144\sqrt{2} \left(\sqrt{1 - \frac{\sqrt{3}}{2}} + \frac{\sqrt{2}}{2} + 1 + \frac{\sqrt{6}}{2} + \sqrt{1 + \frac{\sqrt{3}}{2}}\right) + 144$

 
To simplify the two nested radicals, add them, and call the sum $x$:

 $x = \sqrt{1 - \frac{\sqrt{3}}{2}} + \sqrt{1 + \frac{\sqrt{3}}{2}}$

 Squaring both sides, the F and L part of FOIL causes the radicals to cancel, leaving:

 $x^2 = 2 + 2\sqrt{\left(1 - \frac{\sqrt{3}}{2}\right)\left(1 + \frac{\sqrt{3}}{2}\right)}$

 $x^2 = 2 + 2\sqrt{1 - \frac{3}{4}}$

 $x^2 = 2 + 2\sqrt{\frac{1}{4}}$

 $x^2 = 2 + 2\cdot\frac{1}{2}$

 $x = \sqrt{3}.$

 
Plugging that sum $x$ back into the equation for $l$, we find

 $l = 144\sqrt{2} \left(\frac{\sqrt{2}}{2} + 1 + \frac{\sqrt{6}}{2} + \sqrt{3}\right) + 144$

 $l = 144 + 144\sqrt{2} + 144\sqrt{3} + 144\sqrt{6} + 144$.

 
Thus, the desired quantity is $144\cdot 5 = \boxed{720}$.

 





\textbf{Solution 3}

Begin as in solution 2, drawing a triangle connecting the center of the 12-gon to two vertices of the 12-gon. Apply law of cosines on $\theta = {30^{\circ}, 60^{\circ}, 90^{\circ}, 120^{\circ}, 150^{\circ}, 180^{\circ}}$ to get $d^2 = 288 - 288 \cos \theta$ where $d$ is the diagonal or sidelength distance between two points on the 12-gon. Now, $d=\sqrt{288-288 \cos \theta}$. Instead of factoring out $288$ as in solution 2, factor out $\sqrt{576}$ instead-the motivation for this is to make the expression look like the half angle identity, and the fact that $\sqrt{576}$ is an integer doesn't hurt. Now, we have that $d=24 \sin \frac{\theta}{2}$, which simplifies things quite nicely. Continue as in solution 2, computing half-angle sines instead of nested radicals, to obtain $\boxed{720}$.

 





\textbf{Solution 4 (Another Way to Start Solution 1)}

Plotting this regular $12$-gon on the complex plane with center as origin, and a vertex on the $x$-axis. We have that the vertices are the of the form $12e^{\frac{2\pi k}{12} i}.$ All of the distinct diagonals can be represented with one vertex as $12e^{\frac{2\pi 0}{12} i} = 12.$ Then the each diagonal can be expressed as, 
 \begin{align*} 12 \sqrt{(\cos{0} - \cos{\frac{2\pi k}{12}})^2 + (\sin{\frac{2\pi k}{12}} - \sin{0})^2} &=  12 \sqrt{(\cos{0})^2 + (\cos{\frac{2\pi k}{12}})^2 + (\sin{\frac{2\pi k}{12}})^2 + (\sin{0})^2 - 2(\cos{\frac{2\pi k}{12}}\cos{0} + \sin{\frac{2\pi k}{12}}\sin{0})} \\ &= 12 \sqrt{2 - 2\cos{\frac{2\pi k}{12}}} \\ &= 12 \sqrt{4\sin^2{\frac{\pi k}{12}}} \\ &= 24\sin{\frac{\pi k}{12}}.\end{align*}
 The rest follows as Solution 1.

 




	\item (\textit{AIME I 1990}) Let $T = \{9^k : k ~ \mbox{is an integer}, 0 \le k \le 4000\}$.  Given that $9^{4000}_{}$ has 3817 digits and that its first (leftmost) digit is 9, how many elements of $T_{}^{}$ have 9 as their leftmost digit?


 
 \textbf{Answer: }184



\textbf{Solution 1}

\textbf{Lemma}: For all positive integers n, there's exactly one n-digit power of 9 that does not have a left-most digit 9

 \textit{(Not-so-rigorous) Proof: One can prove by contradiction that there must be at least either one or two n-digit power of 9 for all n.}If there is exactly 1 n-digit power of 9, then such a number $m$ cannot begin with 9 since that would result in $\frac{m}{9}$ also being an n-digits, hence a contradiction. Therefore, this single n-digit power of 9 must not begin with 9. In the case that there are 2 n-digit powers of 9, we have already discovered that one of them must begin with 9, let it be $9^k = m$. The integer $9^{k - 1} = \frac{m}{9}$ must then be an n-digit number that begins with $1$. \textit{}

 
Using this lemma, out of the 4001 powers of 9 in T, exactly 3817 must not begin with 9 (one for each digit). Therefore, the answer is $4001 - 3817 = \boxed{184}$ numbers have 9 as their leftmost digits.

 





\textbf{Solution 2}

Let's divide all elements of $T$ into sections. Each section ranges from $10^{i}$ to $10^{i+1}$ And, each section must have 1 or 2 elements. So, let's consider both cases. 

 If a section has 1 element, we claim that the number doesn't have 9 as the leftmost digit. Let this element be $9^{k}$ and the section ranges from $10^{i}$ to $10^{i+1}$. To the contrary, let's assume the number ($9^{k}$) does have 9 as the leftmost digit. Thus, $9 \cdot 10^i \leq 9^k$. But, if you divide both sides by 9, you get $10^i \leq 9^{k-1}$, and because $9^{k-1} < 9^{k} < 10^{i+1}$, so we have another number ($9^{k-1}$) in the same section ($10^i \leq 9^{k-1} < 10^{i+1}$). Which is a contradiction to our assumption that the section only has 1 element. So in this case, the number doesn't have 9 as the leftmost digit.

 If a section has 2 elements, we claim one has to have a 9 as the leftmost digit, one doesn't. Let the elements be $9^{k}$ and $9^{k+1}$, and the section ranges from $10^{i}$ to $10^{i+1}$.So We know $10^{i} \leq 9^{k} < 9^{k+1} < 10^{i+1}$. From $10^{i} \leq 9^{k}$, we know $9 \cdot 10^{i} \leq 9^{k+1}$, and since $9^{k+1} < 10^{i+1}$. The number($9^{k+1}$)'s leftmost digit must be 9, and the other number($9^{k}$)'s leftmost digit is 1.

 There are total 4001 elements in $T$ and 3817 sections that have 1 or 2 elements. And, no matter how many elements a section has, each section contains exactly one element that doesn't begin with 9. We can take 4001 elements, subtract 3817 elements that don't have 9 as the leftmost digit, and get $\boxed{184}$ numbers that have 9 as the leftmost digit.

 - AlexLikeMath

 





\textbf{Solution 3}

We know that $9$ is very close to $10$. But we also know that for all $k \in \mathbb{Z}^+$, $10^k$ has $k+1$ digits. Hence $10^{4000}$ has $4001$ digits. Now, notice that if a power of $9^n$ has $9$ as the leftmost digit, then $9^{n-1}$ has $1$ as the leftmost digit and will preserve the same number of digits (just due to division). Hence we find that consecutive powers of $9$ are "supposed" to increase by one digit, however, when consecutive powers of $9$ have leading digit $1$ then leading digit $9$ they "lose" this one digit. Therefore, it suffices to find the number of times that $9^k$ has "lost" a digit since the number of times that this occurs shows us the number of pairs where the leading digit of $9^k$ is $1$ and the leading digit of $9^{k+1}$ is $9$. Hence, we find that $9^{4000}$ is "supposed" to have $4001$ digits, but only has $3817$ digits. Thus, the number of times $9^k$ has lost a digit is $4001-3817 = 184$. Thus, $\boxed{184}$ elements of $T$ have $9$ as their leading digit.

 




	\item (\textit{AIME I 1990}) The rectangle $ABCD^{}_{}$ below has dimensions $AB^{}_{} = 12 \sqrt{3}$ and $BC^{}_{} = 13 \sqrt{3}$.  Diagonals $\overline{AC}$ and $\overline{BD}$ intersect at $P^{}_{}$.  If triangle $ABP^{}_{}$ is cut out and removed, edges $\overline{AP}$ and $\overline{BP}$ are joined, and the figure is then creased along segments $\overline{CP}$ and $\overline{DP}$, we obtain a triangular pyramid, all four of whose faces are isosceles triangles.  Find the volume of this pyramid.


 
                    \begin{figure}[H]
                        \centering
                        \includegraphics[width=0.5\textwidth]{Problem/AIME_1990_I_14_Problem_Figure_1.png}
                    \end{figure}
                    


 
 \textbf{Answer: }594



\textbf{Solution 1(Synthetic)}


                    \begin{figure}[H]
                        \centering
                        \includegraphics[width=0.5\textwidth]{Solution/AIME_1990_I_14_Solution_Figure_1.png}
                    \end{figure}
                    

 Our triangular pyramid has base $12\sqrt{3} - 13\sqrt{3} - 13\sqrt{3} \triangle$. The area of this isosceles triangle is easy to find by $[ACD] = \frac{1}{2}bh$, where we can find $h_{ACD}$ to be $\sqrt{399}$ by the \href{https://artofproblemsolving.com/wiki/index.php?title=Pythagorean\_Theorem}{Pythagorean Theorem}. Thus $A = \frac 12(12\sqrt{3})\sqrt{399} = 18\sqrt{133}$.

 
                    \begin{figure}[H]
                        \centering
                        \includegraphics[width=0.5\textwidth]{Solution/AIME_1990_I_14_Solution_Figure_1.png}
                    \end{figure}
                     

 To find the volume, we want to use the equation $\frac 13Bh = 6\sqrt{133}h$, so we need to find the height of the \href{https://artofproblemsolving.com/wiki/index.php?title=Tetrahedron}{tetrahedron}. By the Pythagorean Theorem, $AP = CP = DP = \frac{\sqrt{939}}{2}$. If we let $P$ be the center of a \href{https://artofproblemsolving.com/wiki/index.php?title=Sphere}{sphere} with radius $\frac{\sqrt{939}}{2}$, then $A,C,D$ lie on the sphere. The cross section of the sphere that contains $A,C,D$ is a circle, and the center of that circle is the foot of the \href{https://artofproblemsolving.com/wiki/index.php?title=Perpendicular}{perpendicular} from the center of the sphere. Hence the foot of the height we want to find occurs at the \href{https://artofproblemsolving.com/wiki/index.php?title=Circumcenter}{circumcenter} of $\triangle ACD$. 

 From here we just need to perform some brutish calculations. Using the formula $A = 18\sqrt{133} = \frac{abc}{4R}$ (where $R$ is the \href{https://artofproblemsolving.com/wiki/index.php?title=Circumradius}{circumradius}), we find $R = \frac{12\sqrt{3} \cdot (13\sqrt{3})^2}{4\cdot 18\sqrt{133}} = \frac{13^2\sqrt{3}}{2\sqrt{133}}$ (there are slightly \href{https://artofproblemsolving.com/wiki/index.php?title=Law\_of\_Sines}{simpler ways} to calculate $R$ since we have an isosceles triangle). By the Pythagorean Theorem, 

 
 \begin{align*}h^2 &= PA^2 - R^2 \\ &= \left(\frac{\sqrt{939}}{2}\right)^2 - \left(\frac{13^2\sqrt{3}}{2\sqrt{133}}\right)^2\\ &= \frac{939 \cdot 133 - 13^4 \cdot 3}{4 \cdot 133} = \frac{13068 \cdot 3}{4 \cdot 133} = \frac{99^2}{133}\\ h &= \frac{99}{\sqrt{133}} \end{align*}


 Finally, we substitute $h$ into the volume equation to find $V = 6\sqrt{133}\left(\frac{99}{\sqrt{133}}\right) = \boxed{594}$.

 





\textbf{Solution 2}

Let $\triangle{ABC}$ (or the triangle with sides $12\sqrt {3}$, $13\sqrt {3}$, $13\sqrt {3}$) be the base of our tetrahedron. We set points $C$ and $D$ as $(6\sqrt {3}, 0, 0)$ and $( - 6\sqrt {3}, 0, 0)$, respectively. Using Pythagoras, we find $A$ as $(0, \sqrt {399}, 0)$. We know that the \href{https://artofproblemsolving.com/wiki/index.php?title=Vertex}{vertex} of the tetrahedron ($P$) has to be of the form $(x, y, z)$, where $z$ is the \href{https://artofproblemsolving.com/wiki/index.php?title=Altitude}{altitude} of the tetrahedron. Since the distance from $P$ to points $A$, $B$, and $C$ is $\frac {\sqrt {939}}{2}$, we can write three equations using the \href{https://artofproblemsolving.com/wiki/index.php?title=Distance\_formula}{distance formula}:

 
 \begin{align*} x^{2} + (y - \sqrt {399})^{2} + z^{2} &= \frac {939}{4}\\ (x - 6\sqrt {3})^{2} + y^{2} + z^{2} &= \frac {939}{4}\\ (x + 6\sqrt {3})^{2} + y^{2} + z^{2} &= \frac {939}{4} \end{align*}


 Subtracting the last two equations, we get $x = 0$. Solving for $y,z$ with a bit of effort, we eventually get $x = 0$, $y = \frac {291}{2\sqrt {399}}$, $z = \frac {99}{\sqrt {133}}$.Since the area of a triangle is $\frac {1}{2}\cdot bh$, we have the base area as $18\sqrt {133}$. Thus, the volume is $V = \frac {1}{3}\cdot18\sqrt {133}\cdot\frac {99}{\sqrt {133}} = 6\cdot99 = 594$.

 





\textbf{Solution 3 (Law of Cosines)}

Let $X$ be the apex of the pyramid and $M$ be the midpoint of $\overline{CD}$. We find the side lengths of $\triangle XMP$.

 $MP = \frac{13\sqrt3}{2}$. $PX$ is half of $AC$, which is $\frac{\sqrt{3\cdot13^2+3\cdot12^2}}{2} = \frac{\sqrt{939}}{2}$. To find $MX$, consider right triangle $XMD$; since $XD=13\sqrt3$ and $MD=6\sqrt3$, we have $MX=\sqrt{3\cdot13^2-3\cdot6^2} = \sqrt{399}$.

 Let $\theta=\angle XPM$. For calculating trig, let us double all sides of $\triangle XMP$. By Law of Cosines, $\cos\theta = \frac{3\cdot13^2+939-4\cdot399}{2\cdot13\cdot3\sqrt{313}}=\frac{-150}{6\cdot13\sqrt{313}}=-\frac{25}{13\sqrt{313}}$. 

 Hence, 
 \[\sin\theta = \sqrt{1-\cos^2\theta}\]
 \[=\sqrt{1-\frac{625}{13^2\cdot313}}\]
 \[=\frac{\sqrt{13^2\cdot313-625}}{13\sqrt{313}}\]
 \[=\frac{\sqrt{3\cdot132^2}}{13\sqrt{313}}\]
 \[=\frac{132\sqrt3}{13\sqrt{313}}\]Thus, the height of the pyramid is $PX\sin\theta=\frac{\sqrt{939}}{2}\cdot\frac{132\sqrt3}{13\sqrt{313}}=\frac{66\cdot3}{13}$. Since $[CPD]=3\sqrt{3}\cdot13\sqrt{3}=9\cdot13$, the volume of the pyramid is $\frac13 \cdot 9\cdot13\cdot \frac{66\cdot3}{13}=\boxed{594}$.

 




	\item (\textit{AIME I 1990}) Find $ax^5 + by^5$ if the real numbers $a,b,x,$ and $y$ satisfy the equations

 \begin{align*} ax + by &= 3, \\ ax^2 + by^2 &= 7, \\ ax^3 + by^3 &= 16, \\ ax^4 + by^4 &= 42. \end{align*}



 
 \textbf{Answer: }020



\textbf{Solution 1}

Set $S = (x + y)$ and $P = xy$. Then the relationship

 
 \[(ax^n + by^n)(x + y) = (ax^{n + 1} + by^{n + 1}) + (xy)(ax^{n - 1} + by^{n - 1})\]

 can be exploited:

 
 \begin{eqnarray*}(ax^2 + by^2)(x + y) & = & (ax^3 + by^3) + (xy)(ax + by) \\ (ax^3 + by^3)(x + y) & = & (ax^4 + by^4) + (xy)(ax^2 + by^2)\end{eqnarray*}

 Therefore:

 
 \begin{eqnarray*}7S & = & 16 + 3P \\ 16S & = & 42 + 7P\end{eqnarray*}

 Consequently, $S = - 14$ and $P = - 38$. Finally:

 
 \begin{eqnarray*}(ax^4 + by^4)(x + y) & = & (ax^5 + by^5) + (xy)(ax^3 + by^3) \\ (42)(S) & = & (ax^5 + by^5) + (P)(16) \\ (42)( - 14) & = & (ax^5 + by^5) + ( - 38)(16) \\ ax^5 + by^5 & = & \boxed{020}\end{eqnarray*}

 





\textbf{Solution 2}

A \href{https://artofproblemsolving.com/wiki/index.php?title=Linear\_recurrence}{ recurrence} of the form $T_n=AT_{n-1}+BT_{n-2}$ will have the closed form $T_n=ax^n+by^n$, where $x,y$ are the values of the starting term that make the sequence geometric, and $a,b$ are the appropriately chosen constants such that those special starting terms linearly combine to form the actual starting terms and match the original equations.

 Suppose we have such a recurrence with $T_1=3$ and $T_2=7$. Then $T_3=ax^3+by^3=16=7A+3B$, and $T_4=ax^4+by^4=42=16A+7B$. 

 Solving these simultaneous equations for $A$ and $B$, we see that $A=-14$ and $B=38$. So, $ax^5+by^5=T_5=-14(42)+38(16)= \boxed{020}$.

 Edits-Amon26

 





\textbf{Solution 3}

Using factoring formulas, the terms can be grouped. First take the first three terms and sum them, getting:

 $a(x^3 + x^2 + x) + b(y^3 + y^2 + y) = 16$

 $ax\left( \frac{x^3-1}{x-1} \right) + by\left( \frac{y^3-1}{y-1} \right) = 16$.

 Similarly take the first two terms, yielding:

 $ax \left( \frac{x^2-1}{x-1} \right) + by \left( \frac{y^2-1}{y-1} \right) = 10$.

 Lastly take an alternating three-term sum,

 $a(x^3 - x^2 + x) + b(y^3 - y^2 + y) = 12$$ax \left( \frac{x^3+1}{x+1} \right) + by \left( \frac{y^3+1}{y+1} \right) = 12$.

 Now to get the solution, let the answer be $k$, so 

 $ax \left( \frac{x^4-1}{x-1} \right) + by \left(\frac{y^4-1}{y-1} \right) = 68$.

 Multiplying out this expression gets the answer term and then a lot of other expressions, which can be eliminatedas done in the first solution.

 





\textbf{Solution 4}

We first let the answer to this problem be $k.$ Multiplying the first equation by $x$ gives $ax^2 + bxy=3x$.

 Subtracting this equation from the second equation gives $by^2-bxy=7-3x$. Similarly, doing the same for the other equations, we obtain:$by^2-bxy=7-3x$, $by^3-bxy^2=16-7x$, $by^4-bxy^3=42-16x$, and $by^5-bxy^4=k-42x$

 
Now lets take the first equation. Multiplying this by $y$ and subtracting this from the second gives us  $by^3-bxy^2=(7-3x)y$. We can also obtain  $by^4-bxy^3=(16-7x)y$. 

 Now we can solve for $x$ and $y$! $(7-3x)y = 16-7x$ and $(16-7x)y=42-16x$. Solving for $x$ and $y$ gives us $(-7+\sqrt{87},-7-\sqrt{87})$ (It can be switched, but since the given equations are symmetric, it doesn't matter). $k-42x= (42-16x)y$, and solving for $k$ gives us $k= \boxed{020}$.

 




\end{enumerate}